\documentclass[14pt]{extarticle}
\usepackage{fontspec}
\setmainfont{Times New Roman}
\usepackage[a4paper, top=1cm, left=1cm, right=1cm, bottom=2cm]{geometry}
\usepackage{lib}
\madethi{000}
\begin{document}
\begin{khuvuc}{Bài tập}
    \CauTuLuan{1}{
        \item Viết pt1
    }
    {Lời giải}
    \CauTuLuan{2}{
        \item Viết pt1
    }
    {Lời giải}
    \CauTuLuan{
        Cho hình chóp $S.ABCD$ có đáy ABCD là hình chữ nhật $AB=a, AD=a\sqrt{3}$. Cạnh bên $SA$ vuông góc với đáy và $SA=a$. Hãy tính:
    }{
        \item Góc giữa hai mặt phẳng $(SCD)$ và $(ABCD)$.
        \item Góc giữa hai mặt phẳng $(SBC)$ và $(SAD)$.
    }
    {Lời giải}
    \CauTuLuan{
        Cho hình chóp $S.ABCD$ có đáy $ABCD$ là hình thang vuông tại $A, AB=BC=a;AD=2AB$
        và hai mặt bên $(SAB),(SAD)$ vuông góc với đáy và $SA=a\sqrt{2}$:
    }{
        \item Tính tang của góc $\varphi$ giữa $(SBC)$ và $(ABCD)$.
    }{
        Lời giải
    }
    \CauTuLuan{
        Cho tứ diện $S.ABC$ có các cạnh $SA, SB, SC$ đôi một vuông góc và $SA=SB=SC=1$.
    }{
        \item Tính cosin của góc giữa hai mặt phẳng $(SAB)$ và $(ABC)$.
    }{
        Lời giải
    }
    \CauTuLuan{
        Cho hình chóp $S.ABCD$ có đáy $ABCD$ là hình vuông tâm $O$, cạnh a.
        Đường thẳng $SO$ vuông góc với mặt phẳng đáy và $SO=\dfrac{a\sqrt{3}}{2}$.
    }{
        \item Tính tang của góc $\varphi$ giữa $(SBC)$ và $(ABCD)$.
    }{
        Lời giải
    }
    \CauTuLuan{
        Cho lăng trụ đứng $ABC.A'B'C'$, đáy $ABC$ là tam giác cân $AB=AC=a, BAC=12^\circ, BB'=a, I$ là trung điểm
        của $CC'$.
    }{
        \item Tính cosin của góc giữa hai mặt phẳng $(ABC)$ và $(A'B'I')$.
    }{
        Lời giải
    }
    \CauTuLuan{
        Cho hình lập phương $ABCD.A'B'C'D'$ có cạnh bằng $a$.
    }{
        \item Tính số đo góc giữa $(BA'C)$ và $DA'C)$
    }{
        Lời giải
    }
\end{khuvuc}
\end{document}